\frontsection{Введение}
Практикум подготовлен по дисциплине «Прикладной искусственный интеллект». Искусственный интеллект объединяет методы решения задач, для которых нужны распознавание, прогнозирование или принятие решений. Машинное обучение --- одно из его направлений: модель находит закономерности по примерам. Три работы этого практикума посвящены именно машинному обучению и нейронным сетям на Python.

Практикум учит проходить полный цикл работы: понять данные, определить проверяемую задачу, построить модель, провести сопоставимые эксперименты и объяснить ошибки. Рабочий код служит инструментом исследования. Итог работы --- вывод, который можно проверить по сохранённым данным и результатам.

Три работы связаны общей логикой, но используют разные исходные данные и разные виды ответа. В первой модели получают уже вычисленные признаки сигналов и возвращают один из шести классов активности. Во второй сеть предсказывает непрерывную цену. В третьей признаки изображения формируются внутри свёрточной сети, а ответом служит вид цветка.

\begin{table}[H]\centering\small
\caption{Карта практикума}
\begin{tabularx}{\textwidth}{@{}>{\raggedright\arraybackslash}p{0.07\textwidth}>{\raggedright\arraybackslash}p{0.21\textwidth}Y Y@{}}
\toprule
№ & Задача & Главное новое умение & Проверяемый результат \\
\midrule
1 & Классификация активности & Сравнение алгоритмов; разделение по участникам & Отчёт по шести классам и матрица ошибок \\
2 & Регрессия стоимости & Подготовка смешанных признаков; обучение полносвязной сети & MAE/RMSE, анализ ошибок, таблица прогнозов \\
3 & Классификация цветов & Свёртки, перенос обучения, управление режимами модели & Сравнение трёх архитектур на одной валидации \\
\bottomrule
\end{tabularx}\end{table}

\subsection*{Входная подготовка}
Предшествующий курс машинного обучения не требуется. Нужны функции и коллекции Python, чтение таблиц в pandas, простые операции с массивами NumPy, среднее и представление о производной. Вектор можно сначала понимать как список чисел, матрицу --- как прямоугольную таблицу. Смысл новых понятий объясняется перед их использованием; формулы помогают проследить вычисление.

Начните с раздела «Первое знакомство с машинным обучением». Работы выполняйте по порядку: первая вводит сравнение моделей, вторая --- обучение нейросети, третья --- работу с изображениями. Обзор TabM и TabPFN-3 в конце практикума предназначен для дополнительного чтения.

\subsection*{Общий порядок работы}
\begin{enumerate}
\item Изучите исходные файлы и смысл меток. Проверьте размеры, пропуски, классы и идентификаторы.
\item Зафиксируйте разбиение, метрики, конфигурации и бюджет до сравнения результатов.
\item Выполните простой исходный вариант и проверьте контрольные точки.
\item Меняйте заранее выбранный фактор, сохраняйте результаты каждого опыта.
\item Объясните характерные ошибки и сформулируйте границы вывода.
\end{enumerate}

\begin{keybox}{Как читать примеры кода}
Сначала назовите вход и ожидаемый результат ячейки, затем выполните её и объясните полученное. Учебный модуль --- обычный файл Python с функциями: его можно открыть и прочитать. На первом проходе важнее понять роль функции; затем изучите нужные строки её реализации. Самостоятельная часть --- выбор, изменение модели и объяснение наблюдений.
\end{keybox}

\subsection*{Роль обучающей, валидационной и тестовой частей}
Обучающая часть определяет параметры модели и статистики преобразований. Валидация помогает выбрать настройки. Независимый размеченный тест используют после завершения выбора. Название файла само по себе не определяет его роль: в работе №\,1 \file{test.csv} содержит правильные ответы, а в работе №\,2 --- только признаки новых домов.

Подбор по одной и той же валидации постепенно приспосабливает к ней эксперимент. Поэтому в работах №\,2 и №\,3 результаты на валидации описывают качество выбора в данном опыте; они не объявляются независимой итоговой оценкой. В работе №\,1 предусмотрен отдельный финальный тест.
